\chapter{威胁模型与问题建模}

\section{威胁模型}

本文面向 AMD SEV-SNP 机密虚拟机环境，研究由底层共享硬件资源引发的缓存侧信道问题。系统涉及四类参与实体：运行敏感负载的受害者机密虚拟机（Victim CVM）、发起观测推断的攻击者实体、负责资源管理与调度的宿主机操作系统与 Hypervisor，以及提供缓存层次与一致性维护机制的底层硬件平台。机密计算的核心安全假设之一是将宿主机系统软件排除在可信计算基（TCB）之外，仅将处理器硬件及其安全机制作为信任根\cite{ccc_whitepaper,amd_sev_snp_whitepaper}。本文的威胁模型遵循这一假设，将宿主机操作系统与 Hypervisor 均视为潜在不可信实体。

在该模型下，受害者为机密虚拟机内部的目标程序或其关键执行阶段，攻击者位于宿主侧，在无法直接读取来宾明文内存的前提下，试图通过观测底层微架构副作用推断受害者的访存活动。本文研究的核心不在于对 SEV-SNP 内存加密机制的直接破坏，而在于：当宿主侧实体与机密虚拟机共享处理器缓存、缓存一致性通路以及内存层次结构时，机密虚拟机的访存行为是否会在微架构层面产生可被外部观测者利用的时序差异\cite{chiang2025reloadreload,giner2025coherereload}。

综上，本文的威胁模型可概括为：攻击者在密码学层面无法获取机密虚拟机的明文数据，但可在宿主侧对相关共享硬件资源进行持续采样，通过统计受害者运行期间诱发的缓存一致性变化与访存延迟实施侧信道分析。后续三节将从攻击者能力假设、安全边界划分与攻击路径可行性三个层面对该威胁模型展开论证。

\section{攻击者能力假设}

本文采用能力较强但在现实场景中可成立的攻击者模型。攻击者位于宿主机侧，能够控制宿主机操作系统及 Hypervisor 的运行逻辑，并可部署和执行自定义的用户态或内核态探测程序。这一条件对应内部威胁、受攻陷宿主机管理层或共谋云基础设施等实际场景，符合机密计算威胁模型对宿主机不可信的基本假设\cite{kaplan2016amd,amd_sev_snp_whitepaper}。

在资源可见性方面，攻击者能够获取目标机密虚拟机相关内存页在宿主侧的地址映射关系，并据此建立针对目标页的受控访问路径。需要说明的是，敏感信息所在页面的精确定位不属于本文的核心研究问题，相关工作（如 SEV-Step\cite{wilke2024sevstep}）已给出页面定位流程，本文后续也将给出一种基于粗粒度枚举的敏感页面定位方法。在此基础上，攻击者可沿受控路径执行重复探测，测量访问延迟、捕捉缓存冲突行为，并记录由一致性维护触发的时序变化。进一步地，攻击者可通过控制探测线程的处理器亲和性、采样频率与测量窗口，降低系统调度带来的随机噪声，从而提升信号的稳定性与可区分性。对于目标机密虚拟机的调度时机，攻击者还可通过 vCPU 固定与资源竞争等手段施加一定影响，以缩短侧信道观测窗口并提高攻击时效性。

本文不考虑以下超出上述范围的能力：攻击者不具备破坏 SEV-SNP 内存加密引擎或 RMP 校验机制的能力，不能通过密码学攻击或固件漏洞直接恢复来宾明文数据，也不涉及物理探针、故障注入或其他带外硬件攻击手段。本文所分析的攻击路径，本质上是在不突破机密计算硬件保护边界的前提下，利用共享微架构资源的可观测性实施侧信道推断，而非对机密虚拟机保护机制的直接失效加以利用。

\section{安全边界分析}

安全边界分析的目的在于分析机密虚拟机防御和攻击的能力边界：一类是 SEV-SNP 硬件机制所能有效阻断的攻击路径，另一类是机制设计中未予覆盖、攻击者在现有保护框架下仍可利用的信息通道。

\subsection{SEV-SNP 的保护边界}

SEV-SNP 为机密虚拟机提供了多层次的硬件安全保障。在内存机密性方面，内存控制器中集成的 AES 加密引擎为每个机密虚拟机分配独立密钥，物理内存中存储的始终是密文；即使攻击者以宿主机最高权限直接读取特定物理地址，也只能获得无意义的密文内容\cite{kaplan2016amd,amd_sev_snp_whitepaper}。在内存完整性方面，反向映射表（RMP）记录每个物理页的归属状态与访问权限，硬件在地址转换阶段对映射合法性进行校验，从而阻止恶意 Hypervisor 通过页面重映射、别名映射或回放等方式读取或篡改来宾私有页的内容\cite{amd_sev_snp_whitepaper}。在执行状态保护方面，SEV-ES 及 SEV-SNP 对 VMEXIT 过程中的寄存器上下文实施加密保护，限制 Hypervisor 在上下文切换时对来宾 CPU 状态的可见性，并通过 GHCB 机制将必要的信息交换约束在受控接口内\cite{amd_sev_es_whitepaper}。SEV-SNP 在密码学与完整性层面都为机密虚拟机提供了较为完备的隔离保护，攻击者无法通过直接读取物理内存、操纵页表映射或截获寄存器状态等传统特权软件手段获取来宾明文信息。上述防御措施构成了本文威胁模型中攻击者能力的硬性保护上限。

\subsection{保护边界之外的可观测性}

然而，SEV-SNP 的安全设计主要针对的是对内存内容的直接访问与篡改，并未对底层共享硬件资源的时序行为实施隔离，导致以下微架构行为对宿主侧保持可观测的特征。

\textbf{缓存一致性引发的强制逐出效应。}在 SEV-SNP 平台上，同一物理地址对应着两种数据视图：机密虚拟机访问的是该地址的明文缓存副本，而宿主侧实体访问的是对应的密文副本。当两类副本同时存在于缓存中时，硬件需要通过缓存逐出来维持跨加密域的一致性，并以此避免数据视图不一致问题。这种强制逐出行为会使随后对该地址的访问产生显著的延迟惩罚，而该延迟可在宿主侧被直接测量\cite{giner2025coherereload}。换言之，机密虚拟机的访存活动可以通过触发一致性状态变化，在不泄露任何明文内容的情况下，留下可被外部观测的时序痕迹。

\textbf{物理内存竞争引发的访问延迟。}当机密虚拟机与宿主侧实体在同一时间窗口内并发访问同一物理内存块时，内存控制器需通过硬件仲裁协调两次请求，从而在物理链路上引入附加等待延迟。这一竞争效应与缓存加密状态无关，仅取决于物理访问的时序重叠程度，因而不受 SEV-SNP 任何保护机制的约束\cite{chiang2025reloadreload,ge2018survey}。攻击者可以通过反复触发并测量这一竞争延迟，在统计层面区分受害者是否正在访问目标内存区域。

\textbf{缓存占用与踢出模式。}尽管 SEV-SNP 阻止了对缓存内容的直接读取，机密虚拟机的访存活动仍会影响共享 LLC 的占用状态。当受害者以特定模式访问某内存集合时，其对末级缓存缓存集的占用可能导致攻击者预置数据被替换，进而通过测量重载延迟推断受害者的访问集合\cite{osvik2006cache}。这一通道与传统 Prime+Probe 的基本原理一致，不受内存加密影响。

上述效应说明，SEV-SNP 的加密与完整性保护在阻断直接信息读取的同时，并未消除由硬件共享结构产生的间接时序关联。这种“密码学安全、微架构可观测”的不对称性，正是本文研究的核心安全边界所在。表\ref{tab:sev_snp_boundary}对上述边界进行了归纳。

\begin{table}[htbp]
  \centering
  \renewcommand{\arraystretch}{1.2}
  \caption{SEV-SNP 安全机制的保护范围与边界}
  \label{tab:sev_snp_boundary}
  \begin{tabular}{p{3.2cm}p{3.6cm}p{3.6cm}}
    \hline
    攻击路径 & 所需能力 & SEV-SNP 应对机制\\
    \hline
    直接读取物理内存明文
    & 宿主机访问物理地址
    & 内存加密引擎（AES per-VM key）\\

    页表重映射与别名攻击
    & Hypervisor 操控 NPT
    & RMP 完整性校验 \\

    截获 VMEXIT 寄存器状态
    & Hypervisor 读取上下文
    & SEV-ES 寄存器加密 / GHCB \\

    跨域缓存一致性时序观测
    & 宿主侧时延测量
    & 无专门隔离机制 \\

    物理内存竞争时延侧信道
    & 并发访存 + 计时
    & 无专门隔离机制 \\

    LLC 占用状态推断
    & Prime+Probe 类探测
    & 无专门隔离机制 \\
    \hline
  \end{tabular}
\end{table}

\section{攻击可行性分析}

上述对 AMD SEV-SNP 的分析表明，攻击者虽无法突破密码学保护直接读取来宾数据，却可通过宿主侧时延测量持续观测跨加密域微架构行为。因此，本节在此基础上，从信号形成条件、可区分性、先验需求与噪声约束四个维度，对缓存侧信道攻击的实际可行性进行理论分析。

\subsection{信号形成条件}

侧信道攻击的前提是泄露信号的稳定形成。本文关注的信号来源于两类机制，二者在触发条件上有所不同。

对于跨域一致性逐出信号，其形成需要满足以下条件：宿主侧实体与机密虚拟机在时间上交替访问同一物理地址对应的不同加密域副本，且两次访问之间的间隔足够短，使得前一次访问的缓存副本尚未被系统自然淘汰。在 SEV-SNP 架构下，这一条件相对容易成立——攻击者可主动发起探测访问，并将探测时机对齐到受害者 vCPU 调度窗口附近，从而最大化两类副本在缓存中共存的概率。当条件成立时，缓存一致性协议将强制驱逐与探测访问冲突的加密域缓存行，使得攻击者随后的重载访问产生可测量的延迟惩罚\cite{giner2025coherereload}。

对于内存竞争信号，其形成条件是受害者与攻击者在时间上近似同步地访问同一物理内存块。DRAM 内存控制器在处理并发请求时需通过总线仲裁排队，这一物理排队延迟与缓存状态无关，仅依赖物理访问的时序重叠。当重叠发生时，攻击者测得的访问时延将显著高于非竞争基线。当未发生重叠时，时延则接近正常缓存未命中延迟\cite{chiang2025reloadreload}。

\subsection{信号可区分性分析}

为了实现侧信道攻击，攻击者还需要能够在噪声环境中稳定区分访问目标区域和未访问目标区域两种状态所对应的时延分布。从理论上看，上述两类机制均会导致访问延迟分布的显著偏移：跨域一致性逐出效应通常引入数十至数百纳秒的额外惩罚，内存竞争效应则可产生与 DRAM 访问时间量级相当的附加延迟。这些延迟偏移在统计上具有一定幅度，原则上可通过足够数量的重复采样与分布拟合加以检测\cite{chiang2025reloadreload,giner2025coherereload}。

从攻击实现角度出发，攻击者对每次探测结果的判断本质上是一个二元假设检验：设 $H_0$ 表示受害者在当前探测窗口内未访问目标地址（两类惩罚项均为零），$H_1$ 表示受害者正在访问目标地址（跨域一致性逐出惩罚 $\Delta_\text{coh}$ 或总线竞争惩罚 $\Delta_\text{bus}$ 为正值）。攻击者以固定阈值 $\theta$ 对每次观测延迟 $T_\text{obs}$ 进行判断：

\begin{equation}
  \hat{y} =
  \begin{cases} H_1 & T_\text{obs} > \theta \\ H_0 & T_\text{obs} \leq \theta
  \end{cases}
  \label{eq:threshold_detect}
\end{equation}

该框架下，攻击效果由两类指标衡量：检测率 $P_D(\theta) = P(T_\text{obs} > \theta \mid H_1)$ 表示受害者确实访问目标时攻击者正确判断的概率，误报率 $P_{FA}(\theta) = P(T_\text{obs} > \theta \mid H_0)$ 表示受害者未访问时的错误判断概率。能否找到同时满足高 $P_D$ 与低 $P_{FA}$ 的阈值 $\theta$，取决于 $\Delta_\text{coh}$ 与 $\Delta_\text{bus}$ 的叠加幅度相对于系统噪声的信噪比。本文第四章将以此框架为基础，量化不同实验条件下的 $P_D$ 与 $P_{FA}$，评估自适应阈值策略在真实强噪声环境下的实际效果。

然而，实际环境中存在多种噪声来源影响信号可区分性。系统中断、内核调度抢占、其他虚拟机的并发访存以及硬件预取器的干扰，均可能导致时延测量产生偏差。此外，DRAM Interleaving 等机制会将物理连续地址分散映射至不同内存通道或内存设备，使得竞争信号和一致性惩罚的幅度变得极不稳定。这些因素将导致在理想受控环境中可轻易区分的延迟模式，在真实高噪声平台上的信噪比可能显著下降，需要额外的信号增强与分类策略才能维持足够的识别准确率。

\subsection{攻击所需先验条件}

除信号可观测性外，攻击者还需满足一定的先验知识条件方可将侧信道信号转化为有意义的信息推断。

在地址层面，攻击者需要能够将观测到的宿主侧时延变化与特定的来宾内存区域关联起来。如前文所言，攻击者作为 Hypervisor 控制者天然掌握 GPA 到 HPA 映射，进一步将 HPA 与受害者敏感操作绑定，则依赖对目标程序内存布局的了解或粗粒度枚举探测。在访问目标已知（如特定加密库、操作系统内核代码页）的场景下，这一条件较易满足；在目标未知的场景下，则需要额外的定位步骤。

在时序层面，攻击者需要将探测窗口对齐到受害者执行目标操作的时间段附近。通过 vCPU 调度控制或外部触发（如网络请求驱动受害者执行特定路径），攻击者可以在一定程度上控制这一对齐关系。若受害者操作具有可触发性，则时序对齐相对可控；若目标操作在受害者内部自主触发，则攻击者需通过高频轮询或盲窗口覆盖的方式扩大观测范围。

\subsection{可行性分析结论}

综上所述，本文认为该类侧信道攻击在理论上具备可行性，其核心依据可归纳为以下三点：其一，信号的物理成因来自 SEV-SNP 架构中客观存在的跨加密域一致性维护机制与物理内存竞争机制，而非依赖软件漏洞或偶发行为，因此具备可重复触发的结构性基础；其二，两类机制引入的延迟惩罚在幅度上超过正常访存基线，理论上满足统计可分离的基本要求；其三，攻击者的能力假设（Hypervisor 控制权、探测线程部署能力、地址映射可见性）在云威胁模型中具备现实对应。

与此同时，攻击的实际有效性受到真实系统噪声、硬件架构差异与先验条件获取难度的制约。在强噪声真实平台上，原始时延信号的信噪比可能不足以支撑直接阈值判断，需要结合信号过滤、自适应阈值或统计分类等方法加以补偿。因此，本文后续章节将在理论可行性的基础上，进一步通过实验验证信号在真实机密虚拟机环境中的稳定性与实际可利用程度，为攻击能力的评估提供实证支撑。
